\scanpage{27}
Soit $K$ un corps top. isomorphe à $\mathbf C$ [isom non choisi !]

Alors on a le foncteur \inserted{\unsure{mult.}} $\xi=\xi_K$ sur la catégorie des motifs sur $K$, à valeurs dans les $\mathbf Z$-modules, et des isom
\[
 \xi\otimes_{\mathbf Z}\mathbf Z_\ell\overset{\alpha_\ell}{\simeq}\unsure{\xi}_\ell
\]
\[
 \xi\otimes_{\mathbf Z}K\overset{\alpha_\infty}{\simeq}\xi_K
\]
Considérons
\[
 \xi\bigl(H^2(\mathbf P^1,\unsure{\mathbf Z}(\infty))\bigr)=\xi(\mathbf Z(-1))
\]
c'est isomorphe à $\mathbf Z$, l'isom n'étant pas choisi. Le choix d'un tel isom \inserted{``\unsure{au dire des géom.}''} revient à celui d'une \uline{orientation} sur $\mathbf P_K^1$ [qui est de façon naturelle une variété \deleted{algébrique} réelle, isomorphe à $S^2$), ou ce qui revient au $\unsure{=}$, au choix de l'une des deux structures complexes sur $\mathbf P_K^1$, provenant du \inserted{choix d'un} isom $K\simeq\mathbf C$, i.e. le choix d'\uline{une} des racines de $-1$ dans $K$. Comme $\pi\in\mathbf R\subset K$ est bien déterminé comme

\mspage{26}
